\section{Context Free Languages}
We now consider a larger class of languages known as the \textbf{context-free
languages}. These languages can be defined by a recursive notation known as
\textit{context-free grammars}, or a automaton-like notation known as
\textit{pushdown automaton}.\\

Intuitively, the class of CFLs comprises the RLs, and languages in which
symmetry is required, and certain forms of matching (in opposite directions).

\subsection{Context-Free Grammars}
\begin{definition}[CFG]
  A \textit{CFG} $C=\left( V,T,P,S \right) $, where:
  \begin{itemize}
    \item $V$: A finite set of non-terminals
    \item $T$: A finite set of terminals
    \item $P$: A finite set of productions
    \item $S$: A start symbol, $S\in V$
  \end{itemize}
  A production is of the form $A\to\gamma$, where $A\in V\wedge \gamma\in
  {\left( V\cup T \right)} ^{*}$. If we have the productions
  of $A$ $A\to\alpha_1,A\to\alpha_2,\ldots,A\to\alpha_{n}$, we may write
  $A\to\alpha_1|\alpha_2|\ldots|\alpha_{n}$
\end{definition}
\begin{example}[$L(0^{*}1{(0+1)}^{*})$]
  To generate the regular language $L(0^{*}1{(0+1)}^{*})$, consider the CFG $G=\left(
  \left\{ S,A,B \right\},\left\{ 0,1 \right\},P,S \right) $, where $P$
  comprises the productions:
  \[S\to AB,\qquad A\to 0A|\epsilon,\qquad B\to 0B|1B|\epsilon.\] 
\end{example}
\subsubsection{Derivations}
To infer if a certain string is in the language of the CFG, we may apply
\textit{recursive inference} or \textit{derivation}.
\begin{definition}[Recursive Inference]
  We may construct the string recursively from body to head. For example, given
  a production $A\to BC$, we take strings $b\in L(B)$ and $c\in L(C)$ and
  concatenate them in the proper order and infer that the resulting string is
  in $L(A)$.

  \centering
\begin{tabular}{ccccc}
   \toprule
   & String Inferred & In language of & Production used & String(s) used \\
   \midrule
  (i) & $\epsilon$ & $A$ & $A\to\epsilon$ & -- \\
  (ii) & $\epsilon$ & $B$ & $B\to\epsilon$ & -- \\
  (iii) & $0$ & $A$ & $A\to 0A$ & (i) \\
  (iv) & $0$ & $B$ & $B\to 0B$ & (ii) \\
  (v) & $1$ & $B$ & $B\to 1B$ & (ii) \\
  (vi) & $00$ & $A$ & $A\to 0A$ & (iii) \\
  (vii) & $10$ & $B$ & $B\to 1B$ & (iv) \\
  (viii) & $11$ & $B$ & $B\to 1B$ & (v) \\
  (ix) & $11$ & $S$ & $S\to AB$ & (i), (viii) \\
  \vdots & \vdots & \vdots & \vdots & \vdots\\
  \bottomrule
\end{tabular}

\raggedright
\vspace{2mm}
This procedure is known as \textbf{recursive inference}.
\end{definition}

\begin{notation}[$\Rightarrow$] 
  If $A\in V$ such that $A\to \gamma$ is a production of CFG $G$, we denote
  \[\alpha A\beta\Rightarrow_G \alpha\gamma\beta\qquad\forall\alpha,\beta\in
  {\left( V\cup T \right)}^{*}.\]
\end{notation}
\begin{definition}[Derivation]
  We can expand $S$ with one of its productions, and recursively expand each
  non-terminal with the body of its productions until a string of terminals is derived.
  \[S\Rightarrow AB \Rightarrow A1B \Rightarrow A11B \Rightarrow A110B \Rightarrow A110
  \Rightarrow 0A110 \Rightarrow 0110\]
  This is referred to as a \textbf{derivation} of $0110$. We may then define a
  derivation relation $\Rightarrow ^{*}_G$:
  \begin{enumerate}
    \item \textit{Base Case}: $\alpha \Rightarrow _G^{*} \alpha$
    \item \textit{Inductive Step}: If $\alpha \Rightarrow
      _G^{*}\beta\wedge\beta \Rightarrow _G \gamma$, then $\alpha
      \Rightarrow _G^{*}\gamma$.
  \end{enumerate}
  That is, there exists a derivation from $\alpha$ to $\beta$.
\end{definition}

\begin{definition}[Left/Right-most Derivations]
  A \textbf{leftmost derivation} is one in which at each step, the leftmost
  variable is replaced by one of its production bodies. This restricts the
  number of choices we have in deriving a string. Similarly, we can also
  have a \textbf{rightmost derivation}.
\end{definition}

\begin{definition}[Sentential Forms]
  Any string $\alpha\in {\left( V\cup T \right) }^{*}$ such that $S \Rightarrow
  ^{*}\alpha$ is called a \textbf{sentential form}, a string (of terminals or
  non-terminals) derived from the start symbol. We further have the
  analogue for left/right-most derivations, known as a
  \textit{left/right-sentential form}.
\end{definition}

\subsubsection{Parse Trees}
\begin{definition}[Parse Tree]
  We use a tree representation for derivations, where each interior node is
  labelled by a nonterminal in $V$, each leaf is labelled by a symbol in $V\cup
  T$, or $\epsilon$, and the children $X_1,X_2,\ldots,X_k$ (from the left) of
  an interior node $A$ forms the body of the production $(A\to X_1X_2\cdots
  X_k)\in P$. This is known as the \textbf{parse tree}.
  \begin{figure}[h!]
    \centering
    \includegraphics[width=0.17\textwidth]{parse_tree.jpeg}
  \end{figure}
  We call the in-order traversal of its leaves the \textbf{yield} of the parse
  tree.
\end{definition}
Consider the parse trees whose yields are terminal strings, and the root is
labelled by the start symbol. Then the yield is a string in the language of the
grammar.
\begin{remark}
  Derivations are non unique, but parse trees and its equivalent left-most
  derivations, right-most derivations are unique for unambiguous grammars.
\end{remark}

\subsubsection{Language of a CFG}
\begin{definition}[Language of a CFG]
  The \textbf{language} of a CFG $G$ is denoted $L(G)$, and is defined by
  \[L(G)=\left\{ \omega\in T^{*}| S \Rightarrow _G^{*}\omega \right\},\] the
  set of strings that can be derived from the start symbol in finitely many
  steps.
\end{definition}
\begin{remark}
  The language of a CFG is hence the set of all sentential forms that consist
  solely of terminals.
\end{remark}
\begin{definition}[Context Free Language]
  If $L$ is $L(G)$ for some CFG $G$, then $L$ is a \textbf{context-free language}.
\end{definition}

We prove a property of strings in a specified CFL by inducting on the length of
the string to prove membership, and inducting on the number of steps of
derivation to prove the specified property, as illustrated by the following
example.
\begin{example}[Property of Context-Free Language]
  Suppose the CFG $G=\left( \left\{ S \right\}, \left\{ 0,1 \right\},P, S  \right) $, where
$P$ consists the productions \[P\to \epsilon|0|1|0P0|1P1.\] Then $L(G)$ is the
  set of palindromes over $\left\{ 0,1 \right\} $.\\
  \begin{proof}
    ($\implies $) Suppose $\omega$ is a palindrome. We prove by induction on
    $\left| \omega \right| $ that $\omega$ is in $L(G)$.
    \begin{enumerate}
      \item \textit{Basis Step:} If $\left| \omega \right| =0$ or $\left|
        \omega \right| =1$, then $\omega=\epsilon,0,1$. Since we have the
        productions $P\to\epsilon|0|1$, $P \Rightarrow ^{*}$.
      \item \textit{Inductive Step:} Suppose $\left| \omega \right| \ge 2$.
        Since $\omega$ is a palindrome, it begins and end with the same symbol.
        That is, $\omega=0x0$ or $\omega=1x1$, where $x$ is a palindrome over
        $\left\{ 0,1 \right\} $. By inductive hypothesis, $P \Rightarrow ^{*}x$.
        \begin{itemize}
          \item If $\omega=0x 0$, We have $P \Rightarrow 0P0 \Rightarrow ^{*} 0x0=\omega$.
          \item If $\omega=1x1$, We have $P \Rightarrow 1P1 \Rightarrow ^{*} 1x1=\omega$.A
        \end{itemize}
        We have $P \Rightarrow ^{*}\omega$ is in $L(G)$.
    \end{enumerate}
    ($\implies$) Suppose $\omega\in L(G)$. We prove by induction on the number
    of steps in a derivation of $\omega$ that it is a palindrome.
    \begin{enumerate}
      \item \textit{Basis Step:} For one-step derivations, we have $\epsilon,0,1$ which are
        all palindromes.
      \item \textit{Inductive Step:} Suppose that strings $x$ generated by
        $n$-step derivations are palindromes where $n\ge 1$. Consider an
        $(n+1)$-step derivation of $\omega$. Then $\omega=0x0$ or $\omega=1x1$
        since $n+1\ge 2$, and these are the only productions allowing
        additional steps of a derivation. Since $x$ is a palindrome, $\omega$
        is also a palindrome.\qed%
    \end{enumerate}
  \end{proof}
\end{example}
\vspace{-8mm}
\subsubsection{Inference, Derivations, and Parse Trees}
Given a grammar $G=(V,T,P,S)$, the following are equivalent:
 \begin{itemize}
  \item The recursive inference procedure determines $\omega\in T^{*}$ is in the language of
     $A\in V$
   \item $A \Rightarrow _G^{*}\omega$
   \item There exists a left/right-most derivation from $A$ to $\omega$
   \item There exists a parse tree with root $A$ and yield $\omega$
\end{itemize}

\subsection{Ambiguous Grammars}
\begin{definition}[Ambiguous Grammar]
  A grammar $G$ is said to be ambiguous if, there exists at least two different
  parse trees with root labelled $S$ and yield $\omega\in T^{*}$.
  Equivalently, $G$ is ambiguous if there exist at least two left/right-most
  derivations.
\end{definition}

\begin{remark}
  The ambiguity of a grammar is undecidable.  There is no algorithm to decide
  the ambiguity of a CFG, and no general algorithm to remove ambiguity.
  However, there exists certain techniques to eliminate certain kinds of
  ambiguity.
\end{remark}

\subsubsection{Eliminating Ambiguity}
\begin{example}
  Suppose the CFG $G=\left( \left\{ E,I,S \right\} ,\left\{ a,b,0,1 \right\} ,P,S \right) $ where
  $P$ consists the productions
  \[E \to I|E+E|E*E|(E)\qquad I\to a|b|Ia|Ib|I0|I1.\]
  This grammar is ambiguous for two causes: (1) The precedence of operators is
  not defined, and (2) a sequence of identical operators can be grouped from the
  left or the right. To eliminate ambiguity, we impose an order of precedence, and
  left-associativity of operators.

  Precedence can be enforced by introducing different variables,
  each of which represents expressions sharing the same level of ``binding strength'':
  \begin{itemize}
    \item Identifiers and paranthesised expressions are \textit{factors},
      expressions that cannot be broken apart by any adjacent operator.
    \item A product of one or more factors are \textit{terms}, expressions that
      cannot be broken by the $+$ operator.
    \item A sum of one or more terms is then all other possible \textit{expressions} to be
      represented.
  \end{itemize}
  We hence disambiguate with the grammar
  \begin{align*}
    I&\to a|b|Ia|Ib|I0|I1\\
    F&\to I|(E)\\
    T&\to F|T*F\\
    E&\to T|E+T
  \end{align*}
\end{example}

\begin{remark}
  It is difficult to prove disambiguity of a specified grammar.
\end{remark}

\subsubsection{Inherent Ambiguity}
\begin{definition}[Inherently Ambiguous]
  A CFL $L$ is \textbf{inherently ambiguous} if all its grammars are ambiguous.
\end{definition}
\begin{example}[Inherently Ambiguous CFL]
   \[L=\left\{ a^{n}b^{n}c^{m}d^{m}|n\ge 1,m\ge 1 \right\} \cup \left\{ a^{n}b^{m}c^{m}d^{n}|n\ge 1,m\ge 1 \right\}\] 
   is an inherently ambiguous language. 
\end{example}
\begin{remark}
  Proof of inherent ambiguity is complex.
\end{remark}

\subsection{Simplification of CFG}
\subsubsection{Prerequisites}
For a grammar $G=\left( V,T,P,S \right)$, we have the following definitions:
\begin{definition}[Generating Symbols]
  A symbol $A$ is a \textit{generating symbol} if $\exists \omega\in T^{*}$,
  $A\Rightarrow ^{*} \omega$.
\end{definition}
\begin{definition}[Reachable Symbols]
  A symbol $A$ is a \textit{reachable symbol} if $\exists \alpha,\beta\in
  {\left( V\cup T \right)}^{*}$, $S\Rightarrow ^{*} \alpha A\beta$.
\end{definition}
\begin{definition}[Useless Symbols]
  A symbol $a\in V\cup T$ is a \textit{useless symbol} if it does not
  appear in any derivation of a terminal string from the start symbol
  $S\Rightarrow ^{*}\omega$ for any $\omega\in T^{*}$.
\end{definition}
A symbol $A$ is useful (not useless) $\implies$ $A$ is reachable and
generating. The converse is not always true.
\begin{definition}[$\epsilon$-Productions]
  An $\epsilon$-production is a production of the form $A\to\epsilon$ for some
  $A\in V$.
\end{definition}
\begin{definition}[Nullable Symbols]
  A symbol $A$ is a \textit{nullable symbol} if $A\Rightarrow ^{*}\epsilon$.
\end{definition}
\begin{definition}[Unit Productions]
  A unit production is a production of the form $A\to B$ for some
  $A,B\in V$.
\end{definition}
\begin{definition}[Unit Pairs]
  A tuple $\left( A,B \right) $ for some $A,B\in V$ is a unit pair if
  $A\Rightarrow^{*} B$.
\end{definition}

\subsubsection{CFG Simplification Algorithm}
Suppose we are given a grammar $G=\left( V,T,P,S \right)$ such that
$L(G)\neq\emptyset $. Let $G_3=\left( V_3,T_3,P_3,S \right)$ be the grammar obtained by
\textbf{eliminating all $\epsilon$-productions} by first finding all nullable
symbols by iterating the following inductive algorithm:
\begin{enumerate}
  \item \textit{Basis Step:} If $A\to\epsilon$ is a production of $G$, $A$
    is nullable.
  \item \textit{Inductive Step:} If $B\to \beta$ is a production and
    $\beta$ consists only of nullable nonterminals, $B$ is nullable.
\end{enumerate}
For each production $A\to X_1X_2\cdots X_k$, given $m$ of the $k$ $X_i$s
are nullable, create $2^{m}$ versions of the production of all possible
combinations of presence/absence of the nullable $X_i$s. If all symbols are
nullable, the case of $A\to\epsilon$ is not included.
\begin{example}[Removing $\epsilon$-productions]
  Suppose $S\to AB$, where $A,B$ are nullable symbols. In the new grammar,
  we have the productions $S\to AB | A | B$.
\end{example}
Then $G_3$ does not contain any $\epsilon$-productions, and
$L(G_3)=L(G)-\left\{ \epsilon \right\} $.\\

\begin{proof}
  We prove the more general statement that
  \[A\Rightarrow_{G_3}^{*}\omega\quad\iff\quad A\Rightarrow_{G}^{*}\omega\wedge
  \omega\neq\epsilon.\]
  ($\implies$) Suppose $A\Rightarrow_{G_3}^{*}$. $\omega\neq\epsilon$ since
  $G_3$ has no $\epsilon$-productions. Inducting on the number of steps of the
derivation $A\Rightarrow_{G_3}^{*}\omega$:
  \begin{enumerate}
    \item \textit{Basis Step:} There is a production $A\to\omega$ in $G_3$. By
      construction of $G_3$, there is some production $A\to\alpha$ in G such
      that $\alpha$ is $\omega$ with some number of nullable symbols
      interspersed. We have $A\Rightarrow_G\alpha\Rightarrow _G^{*}\omega$ by
      deriving $\epsilon$ from all the nullable symbols.
    \item \textit{Inductive Step:} Suppose a derivation of $n>1$ steps. The derivation
      must look like
      \[A\Rightarrow_{G_3}X_1X_2\cdots X_k\Rightarrow_{G_3}^{*}\omega.\]
      Similarly, this must come from a production $A\to Y_1Y_2\cdots Y_m$ in
      $G$, where the $Y$s are the $X$s in order with some nullable symbols
      interspersed. Let $\omega=\omega_1\omega_2\cdots\omega_k$, where
      $X_i\Rightarrow ^{*}_{G_3}\omega_i$. Each derivation takes fewer than $n$
      steps, and hence $X_i\Rightarrow _G^{*}\omega_i$ by inductive hypothesis.
      \[A\Rightarrow_G Y_1Y_2\cdots Y_m\Rightarrow _G^{*}X_1X_2\cdots
      X_k\Rightarrow ^{*}_G\omega_1\omega_2\cdots\omega_k=\omega.\] 
  \end{enumerate}
  ($\implies$) Suppose $A\Rightarrow ^{*}_G \omega$ and $\omega\neq\epsilon$. Inducting
  on the number of steps of the derivation $A\Rightarrow _{G}^{*}\omega$:
  \begin{enumerate}
    \item \textit{Basis Step:} There is a production $A\to \omega$ in $G $.
      Since $\omega\neq\epsilon$, it is also a production of $G_3$, and
      $A\Rightarrow _{G_3}^{*}\omega$.
    \item \textit{Inductive Step:} Suppose a derivation of $n>1$ steps. The
      derivation must look like  $A\Rightarrow _G Y_1Y_2\cdots
      Y_m\Rightarrow_G^{*}\omega$. We can break up
      $\omega=\omega_1\omega_2\cdots\omega_m$ such that $Y_i\Rightarrow
      _G^{*}\omega_i$. Let $X_1,X_2,\ldots,X_k$ be the $Y_j$s such that
      $w_j\neq\epsilon$. Since $\omega\neq\epsilon$, $k\ge 1$ and $A\to
      X_1X_2\cdots X_k$ is a production in $G_3$. Each derivation
      $Y_j\Rightarrow_G^{*}\omega_j$ takes fewer than $n$ steps, and hence if
      $\omega_j\neq\epsilon$, $Y_j\Rightarrow ^{*}_{G_3}\omega_j$.
      \[A\Rightarrow _{G_3}X_1X_2\cdots X_k\Rightarrow _{G_3}^{*}\omega.\]
  \end{enumerate}
  We have $S\Rightarrow_{G_3}^{*}\omega\iff A\Rightarrow_G^{*}\omega\wedge\omega\neq\epsilon$. Hence $L(G_3)=L(G)-\left\{ \epsilon \right\} $.\qed
\end{proof}\\

Let $G_4=\left( V,T,P_4,S \right)$ be the grammar obtained by
\textbf{eliminating all unit productions} by first finding all unit pairs by
iterating the following inductive algorithm:
\begin{enumerate}
  \item \textit{Basis Step:} $\forall A\in V$, $\left( A,A \right) $ is a unit pair.
  \item \textit{Inductive Step:} If $\left( A,B \right) $ is a unit pair and
    $B\to C$ is a production, for some nonterminal $C$, $\left( A,C \right) $
    is a unit pair.
\end{enumerate}
For each unit pair $\left( A,B \right)$, let $\left\{ A\to\alpha |
B\to\alpha\text{ is a non unit production in $P$} \right\}\subseteq P'$. Then $G'$ does
not contain any unit productions, and $L(G_4)=L(G)$.\\

Let $G_1=\left( V_1,T_1,P_1,S \right) $ be the grammar
obtained with the following steps:
\begin{enumerate}
  \item \textbf{Eliminating non-generating symbols} and all productions
    involving one or more of those symbols. We find all generating symbols by
    iterating the following inductive algorithm:
    \begin{enumerate}
      \item \textit{Basis Step} All terminals $a\in T$ are generating.
      \item \textit{Inductive Step:} If $A\to\alpha$ is a production and
        $\alpha$ consists only of generating symbols, $A$ is generating.
    \end{enumerate}
  \item \textbf{Eliminating non-reachable symbols} in the resultant grammar
    $G_2$. We find all reachable symbols by iterating the following inductive
    algorithm:
    \begin{enumerate}
      \item \textit{Basis Step} $S$ is reachable.
      \item \textit{Inductive Step:} If $A$ is reachable and $A\to\alpha$ is a
        production, every symbol in $\alpha$ is reachable.
    \end{enumerate}
\end{enumerate}
Then $G_1$ has no useless symbols and $L(G_1)=L(G)$.\\

\begin{proof}
  Suppose a symbol in $G_1$, $X\in V_1\cup T_1$. $X$ is generating, and every
  symbol used in the derivation $X\Rightarrow ^{*}\omega$ is also generating.
  We have $X\Rightarrow_{G_2}^{*}\omega$.

  Since $X$ is not eliminated in the second step, $\exists \alpha,\beta\in
  {\left( V\cup T \right)}^{*}$, $S\Rightarrow_{G_2}^{*} \alpha A\beta$, and
  every symbol in the derivation is also reachable, hence
  $S\Rightarrow_{G_1}^{*}\alpha X\beta$.\\

  Since every symbol in $\alpha X\beta$ is reachable and generating (since they
  are in $V_2\cup T_2$)
  \[S\Rightarrow_{G_1}^{*}\alpha X\beta\Rightarrow_{G_1}^{*}xwy\]
  for some $x,\omega,y\in T^{*}$. $X$ is useful in $G_1$.\\

  Additionally, $L( G_1 ) \subseteq L(G)$ since we have only eliminated symbols
  and productions. $L(G) \subseteq L(G_1)$ since every derivation
  $S\Rightarrow_G^{*}\omega$ consists of only useful symbols, hence
  $S\Rightarrow _{G_1}^{*}\omega$. $L(G)=L(G_1)$. \qed
\end{proof}\\

Supposing there is at least one string in $L(G)$, by following this procedure
of (in this order):
\begin{enumerate}
  \item Eliminating $\epsilon$-productions
  \item Eliminating unit productions
  \item Eliminating useless symbols
\end{enumerate}
we obtain a simplified CFG $G'$ with the various simplifications, with
$L(G')=L(G)-\left\{ \epsilon \right\} $.

\begin{remark}
To obtain an equivalent CFG if $\epsilon\in L(G)$, we specify the CFG $G''=\left( V\cup \left\{ S_0 \right\} ,T,P'', S_0 \right) $, where 
\[P''=P\cup \left\{ S_0\to S, S_0\to\epsilon \right\}.\] 
\end{remark}

\subsubsection{Normal Forms}
\begin{definition}[Chomsky Normal Form]
  A CFG $G=\left( V,T,P,S \right) $ is said to be in \textit{Chomsky Normal
  Form} if all productions are of the form 
  \[A\to BC\qquad\text{or}\qquad A\to a\]
  where $A,B,C\in V$ and $a\in T$, and $G$ has no useless symbols.
\end{definition}

CNF guarantees non-decreasing length of the derived string, with respect to the
number of productions used.

\begin{theorem}
  Every nonempty language without $\epsilon$ is $L(G)$ for some CFG $G$ in
  Chomsky Normal Form.
\end{theorem}

Given the simplified CFG without $\epsilon$-productions, unit productions and
useless symbols, we can obtain an equivalent CFG in CNF as follows:
Suppose a production 
\[A\to X_1X_2\cdots X_k.\]
We break the production into the following set of productions:
\begin{gather*}
  A\to Z_1B_2\\
  B_2\to Z_2B_3\\
      \vdots\\
  B_{k-1}\to Z_{k-1}Z_k\\
  Z_i\to X_i,\quad\text{if }X_i\in T\\
  Z_i=X_i,\quad\text{if }X_i\in V
\end{gather*}
with new nonterminals $B_i$ and (possibly) $Z_i$. 

\begin{theorem}[Size of Parse Tree]
  Suppose a parse tree of a Chomsky-Normal-Form Grammar. If the length of the
  longest path from the root to a node is $k$, the size of the string generated
  is at most $2^{k-1}$.
\end{theorem}

\begin{definition}[Greibach Normal Form]
  A CFG $G=\left( V,T,P,S \right) $ is said to be in \textit{Greibach Normal
  Form} if all productions are of the form 
  \[A\to a\alpha.\] 
  where $A\in V$, $a\in T$ and $\alpha\in {\left( V\cup T \right)}^{*}$.
\end{definition}

GNF guarantees a generated terminal in every production used as the leftmost
symbol.

\begin{theorem}
  Every nonempty language without $\epsilon$ is $L(G)$ for some CFG $G$ in
  Greibach Normal Form.
\end{theorem}

% TODO: Add proof of greibach normal form

\subsection{Pushdown Automata}
A pushdown automaton is in essence a NFA with $\epsilon$-transitions and a stack on which
it can store a string of ``stack symbols''.
\begin{definition}[PDA]
  A PDA $P=\left( Q,\Sigma,\Gamma,\delta,q_0,Z_0,F \right) $, where:
  \begin{itemize}
    \item $Q$: A finite set of states
    \item  $\Sigma$: A finite set of input symbols
    \item $\Gamma$: A finite stack alphabet
    \item $\delta$: A transition function taking as argument a triple $\left(
      p,a,X \right)$ representing the current state of the PDA and returns a
      set of pairs $\left( q,\gamma \right) $ of new states and replacement
      stack symbols.  $\delta:Q\times \Sigma  \times \Gamma \to
      \mathcal{P}\left( Q\times \Gamma^{*} \right) $ 
    \item $q_0$: A start state, $q_0\in Q$ 
    \item $Z_0$: A initial symbol on the stack, $Z_0\in \Gamma$ 
    \item $F$: A set of final or accepting states.  $F\in \mathcal{P}\left( Q \right) $
  \end{itemize}
\end{definition}
We have a transition diagram for PDAs, in which:
\begin{itemize}
  \item Nodes correspond to states of the PDA
  \item Doubly circled states are accepting
  \item An arrow labelled ``Start'' indicates the start state
  \item Each arc corresponds to a transition of the PDA:\\
    An arc labelled $a,X/\alpha$ from state $p$ to $q$ indicates $\left(
    p,\alpha \right) \in \delta\left( q,a,X \right) $ 
  \item By convention, we take $Z_0$ to be the start symbol
\end{itemize}
\begin{figure}[h!]
  \centering
  \includegraphics[width=0.5\textwidth]{pda_diag.jpeg}
\end{figure}
\begin{definition}[Instantaneous Descriptions]
  The \textbf{ID} of a PDA is the triple $\left( q,\omega,\gamma \right) $ such
  that $q$ is the current state, $\omega$ the remaining input and $\gamma$ the
  stack contents.
\end{definition}

\begin{notation}[$\vdash$]
  If $\left( p,\alpha \right) \in \delta\left( q,a,X \right)$, we denote
  \[\left( q,a\omega,X,\beta \right) \vdash \left( p,\omega,\alpha\beta
  \right)\qquad \forall \omega\in \Sigma^{*},\beta\in \Gamma^{*}.\]
  This represents one move of the PDA.
\end{notation}
\begin{notation}[$\vdash^{*}_P$]
  To represent zero or moves of the PDA, we further define the relation $\vdash^{*}_P$:
  \begin{enumerate}
    \item \textit{Base Case}: $I\vdash^{*}_P I$ 
    \item \textit{Inductive Step}: For IDs $I,J$, if $\epsilon$ some ID $K$
      such that $I\vdash K$ and $K\vdash^{*}_P J$, then $I\vdash^{*}_P J$
  \end{enumerate}
\end{notation}

\subsubsection{Language of a PDA}
\begin{definition}[Language of a PDA]
  The \textit{language} of a PDA $P$ is denoted $L(P)$ and depends on the modality of acceptance
  of the PDA\@. It is defined by
  \[\left\{ \omega|\left( q_0,\omega,Z_0 \right) \vdash_P^{*}\left( q_f,\epsilon,\alpha \right) \right\},q_f\in F \quad\text{ for acceptance by final state.}.\] 
  \[\left\{ \omega|\left( q_0,\omega,Z_0 \right) \vdash_P^{*}\left( q_f,\epsilon,\epsilon \right) \right\},q_f\in F \quad\text{ for acceptance by empty stack.}.\] 
\end{definition}

\begin{theorem}[Equivalence of Acceptance Modalities]
  The two modalities of acceptance are equivalent, in the sense that $L$ has a
  PDA that accepts it by final state $\iff$ $L$ has a PDA that accepts it by
  empty stack.\\
  Equivalently, the class of languages accepted by final state is equal to the class
  of languages accepted by empty stack.
\end{theorem}
\begin{theorem}
  Suppose a language $L$ has a PDA that accepts it by empty stack. Then $L$ has
  a PDA that accepts it by final state.
\end{theorem}
\begin{proof}
  Suppose $L=L(P_N)$ through acceptance by empty stack for some PDA $P_N=\left(
  Q,\Sigma,\Gamma,\delta_N,q_0,Z_0,F \right) $. Consider the PDA
  \[P_F=\left( Q\cup \left\{ p_0,p_f \right\} ,\Sigma,\Gamma\cup \left\{ X_0
  \right\} ,\delta_F,p_0,X_0,\left\{ p_f \right\}  \right)\]
  such that
  \begin{enumerate}
    \item $P_F$ spontaneously transits to the start state of $P_N$, pushing the start symbol.
      \[\delta_F\left( p_0,\epsilon,X_0 \right) =\left\{ \left( q_0,Z_0 X_0 \right)  \right\}\]
    \item $P_F$ simulates $P_N$.  \[\forall q\in Q,a\in \Sigma\cup \left\{ \epsilon \right\},Y\in
      \Gamma\qquad \delta_N\left( q,a,Y \right) \subseteq \delta_F\left(
    q,a,Y\right) \]
  \item After the stack is empty, transit to the final state of $P_F$.
    \[\forall q\in Q\qquad \left( p_f,\epsilon \right) \in \delta_F\left(
    q,\epsilon,X_0 \right) \]
  \end{enumerate}
  We now show that $w\in L(P_F)\text{ with acceptance by final state } \iff \omega\in L$.\\
  ($\impliedby$) Suppose $\omega\in L$. $\left( q_0,\omega,Z_0 \right) \vdash_{P_N}^{*}$ for
  some state $q$.
  \[\left( q_0,\omega,Z_0X_0 \right) \vdash_{P_N}^{*}\left( q,\epsilon,X_0
  \right).\]
  Since $P_F$ has all the moves of $P_N$, $\left(
  q_0,\omega,Z_0X_0 \right) \vdash_{P_F}^{*}\left(
  q,\epsilon,X_0 \right) $. Hence we have
  \[\left( p_0,\omega,X_0 \right) \vdash_{P_F}\left( q_0,\omega,Z_0X_0 \right) \vdash_{P_F}^{*}\left( q,\epsilon,X_0 \right)\vdash_{P_F}\left( p_f,\epsilon,\epsilon \right)\]
  and $P_F$ accepts $\omega$ by final state.\\
  ($\implies$) Suppose $P_F$ accepts $\omega$ by final state. Since
  transitions to $p_f$ require $X_0$ as the stack symbol, and $X_0$ can only be
  on the bottommost position, any computation of $P_F$ that accepts $\omega$
  must be of the previously shown sequence.  Hence $\omega$ is in $P_N$.
  \begin{figure}[h!]
    \centering
    \includegraphics[width=0.6\textwidth]{pda_empty_final.jpeg}
  \end{figure}
  \qed
\end{proof}

\begin{theorem}
  Suppose a language $L$ has a PDA that accepts it by final state. Then $L$ has
  a PDA that accepts it by empty stack.
\end{theorem}
Suppose $L =L(P_F)$ through acceptance by final state for some PDA
$P_F=\left( Q,\Sigma,\Gamma,\delta_F,q_0,Z_0,F \right) $. Consider the PDA
\[P_N=\left( Q\cup \left\{ p_0,p \right\},\Sigma,\Gamma\cup \left\{ X_0 \right\} ,\delta_N,p_0,X_0 \right)\]
such that
\begin{enumerate}
\item We push the start symbol of $P_F$ onto the stack and going to the start state of $P_F$. \[\delta_N\left( p_0,\epsilon,X_0 \right) =\left\{ q_0,Z_0 X_0 \right\} \]
\item $P_N$ simulates $P_F$. \[\forall q\in Q,a\in \Sigma\cup \left\{ \epsilon \right\},Y\in \Gamma\qquad \delta_F\left( q,a,Y \right) \subseteq \delta_N\left( q,a,Y \right) \]
  \item If $P_F$ accepts, $P_N$ starts emptying its stack without consuming
    input. \[\forall q\in F,Y\in \Gamma\cup \left\{ X_0 \right\}\qquad \left(
  p,\epsilon \right) \in \delta_N\left( q,\epsilon,Y \right) \]
  \item After $P_F$ has accepted, $P_N$ pops every symbol until the stack is
    empty. \[\forall Y\in \Gamma\cup \left\{ X_0 \right\}\qquad \delta_N\left(
    p,\epsilon,Y \right) =\left\{ \left( p,\epsilon \right)  \right\} \]
\end{enumerate}
\begin{figure}[h!]
  \centering
  \includegraphics[width=0.6\textwidth]{pda_final_empty.jpeg}
\end{figure}
It can be shown that $\omega\in L(P_N)$ with acceptance by empty stack $\iff\omega\in L$.
\begin{remark}
  For a given PDA $P$, the languages that $P$ accepts by final state and by
  empty stack are different in general.
\end{remark}


\subsection{CFGs vs PDAs}
\subsubsection{Equivalence of CFGs and PDAs}
\begin{theorem}
  A language $L$ is accepted by some PDA if $L$ is accepted by some CFG\@.
\end{theorem}
\begin{proof}
  Suppose a CFG $G=\left( V,T,G,S \right) $. We construct a PDA $P$ as follows: 
  \[P=\left( \left\{ q \right\} ,T,V\cup T,\delta,q,S \right)\]
  where $\delta$ is defined:
  \begin{enumerate}
    \item For each nonterminal $A\in V$, $\delta\left( q,\epsilon,A \right)
      =\left\{ \left( q,\beta \right) | A\to \beta\in G \right\}$
    \item For each terminal $a\in T$, $\delta\left( q,a,a \right) =\left\{ q,\epsilon \right\} $
  \end{enumerate}
  We prove that $\omega\in L(G)\iff \omega\in L(P)$ with
  acceptance by empty stack.\\

  ($\implies$) Suppose $\omega\in L(G)$. Then $\omega$ has a leftmost derivation
  \[S = \gamma_1 \Rightarrow \gamma_2 \Rightarrow \cdots \Rightarrow \gamma_{n}=\omega.\] 
  We show by induction on $i$ that $\left( q,\omega,S \right) \vdash^{*} \left(
  q,y_i,\alpha_i \right) $, where $\gamma_i=x_i\alpha_i$, $\omega=x_i y_i$,
  $x_i,y_i\in T^{*}$ and $\alpha\in \left( V\cup T \right) ^{*}$. That is,
  $\alpha_i$ is the tail of $\gamma_i$ from the first nonterminal (from the
  left), and $y_i$ is the nonterminal string to be generated from $\alpha_i$.
  \begin{enumerate}
    \item \textit{Basis Step:} $\gamma_1=S$, $x_1=\epsilon$, $y_1=\omega$. By definition,
      $\left( q,\omega,S \right) \vdash^{*}\left( q,\omega,S \right) $.
    \item \textit{Inductive Step:} Suppose by inductive hypothesis that $\left(
      q,\omega,S \right) \vdash^{*}\left( q,y_i,\alpha_i \right) $. $\alpha_i$
      begins with a nonterminal $A$. By replacing $A$ by one of its production
      bodies, say $\beta$ with rule (1) of $\delta$, and removing all terminals at the
      top of the stack with rule (2) of $\delta$, we reach the ID $\left(
      q,y_{i+1},\alpha_{i+1} \right) $.
  \end{enumerate}
  Since $\omega$ is a string of terminals, $\alpha_{n}=\epsilon$, and we have
  $\left( q,\omega,S \right) \vdash^{*}\left( q,\epsilon,\epsilon \right) $.
  $P$ accepts $\omega$ by empty stack.\\

  ($\impliedby$) Suppose $\left( q,x,A \right) \vdash^{*}\left(
  q,\epsilon,\epsilon \right) $. Inducting on the number of moves taken by $P$,
  we have:
  \begin{enumerate}
    \item \textit{Basis Step:} $A\to\epsilon$ is a production of $G$, and the
      move is of rule (1). We have $x=\epsilon$ and $A \Rightarrow \epsilon$.
    \item \textit{Inductive Step:} Suppose $P$ takes $n$ moves. The first move
      must be of rule (1), where $A$ is replaced by one of its production
      bodies. Let the production used be $A\to Y_1 Y_2 \cdots Y_k$ for some
      $Y_i\in \left( V\cup T \right)^{*}$. By taking $x=x_1 x_2\cdots x_k$,
      where each $x_i$ is the portion of the input consumed until $Y_1$ is
      popped off the stack, for all $i=1,2,\ldots,k$:
      \[\left( q,x_i x_{i+1}\cdots x_k,Y_i \right) \vdash^{*}\left(
      q,x_{i+1}\cdots x_k,\epsilon \right).\]
      Each of these sequences are no more than $n-1$ moves (since the PDA
      completes in $n-1$ moves), $Y_i \Rightarrow ^{*}x_i$ for nonterminals
      $Y_i$. For terminals $Y_i$, only one move is involved, matching the one
      symbol against $Y_i$. Hence we have the derivation:
      \[A \Rightarrow Y_1 Y_2\cdots Y_k \Rightarrow ^{*}x_1Y_2\cdots Y_k \Rightarrow ^{*}\cdots \Rightarrow^{*} x_1x_2\cdots x_k.\]
      We have $A \Rightarrow ^{*}x$.
  \end{enumerate}
  By letting $A=S$ and $x=\omega$, since $\omega\in L(P)$ by acceptance by empty stack,
  $S \Rightarrow \omega$ and it is in $L(G)$.\qed
\end{proof}
\begin{theorem}
  A language $L$ is accepted by some CFG if $L$ is accepted by some PDA\@.
\end{theorem}
\begin{proof}
  Suppose a PDA $P=\left( Q,\Sigma,\Gamma,\delta,q_0,Z_0 \right)$ accepting by
  empty stack. We construct a CFG as follows:
  \[G=\left( V,\Sigma,R,S \right)\]
  where $V=\left\{ S \right\} \cup \left\{ \left[ pXq \right] |p,q\in Q,Z\in
  \Gamma \right\} $. Each nonterminal encodes the `net popping' of symbol $X$
  from the stack, with a change in state from $p$ to $q$. The productions of
  $G$ are then:
  \begin{enumerate}
    \item $\forall p\in Q$, $S\to\left[ q_0Z_0p \right] \in R$. That is, we
      generate all strings that pop $Z_0$ while going from state $q_0$ to (any)
      state $p$, $\left( q_0,\omega,Z_0 \right) \vdash^{*}\left(
      p,\epsilon,\epsilon \right) $.
    \item If $\left( r,Y_1Y_2\cdots Y_k \right) \in \delta\left( q,a,X \right) $, where $a\in \Sigma\cup \left\{ \epsilon \right\}, k\in \mathbb{N}$, 
      \[\forall \left( r_1,r_2,\ldots,r_k \right) \in Q^{k},\qquad \left[ qXr_k
      \right] \to a\left[ rY_1r_1 \right] \left[ r_1Y_2r_2 \right]\cdots\left[
    r_{k-1}Y_kr_k \right].\] 
    Intuitively, one of the ways to pop $X$ and go from states $q$ to $r_k$ is
    to read $a$ and then use some input to sequentially pop each of the
    produced symbols.
  \end{enumerate}
  We prove that $\left( q,\omega,X \right) \vdash^{*}\left( p,\epsilon,\epsilon
  \right)\iff\left[ qXp \right] \Rightarrow ^{*}\omega$.\\

  ($\implies$) Suppose $\left( q,\omega,X \right) \vdash^{*}\left(
  p,\epsilon,\epsilon \right)$. Inducting on the number of moves taken by $P$, we have:
  \begin{enumerate}
    \item \textit{Basis Step:} With one step, $\left( p,\epsilon \right) \in \delta\left( q,\omega,X \right)$ and $\omega$ is a single symbol or $\epsilon$. By construction, $\left[ qXp \right] \to \omega$ is a production and $\left[ qXp \right] \Rightarrow \omega$.
    \item \textit{Inductive Step:} Suppose the sequence $\left( q,\omega, X
      \right) \vdash^{*}\left( p,\epsilon,\epsilon \right) $ takes $n>1$ steps.
      The sequence must be of the form
      \[\left( q,\omega, X \right) \vdash \left( r_0,x,Y_1Y_2\cdots Y_k \right)
      \vdash^{*}\left( p,\epsilon, \epsilon \right)\] 
      where $\omega=ax$, $a\in \left\{ \epsilon \right\} \cup \Sigma$. We have 
      $\left( r_0,Y_1Y_2\cdots Y_k \right)\in \delta\left( q,a,X \right) $ and
      by construction, we have the productions $\left[ qXp \right] \to a\left[
      r_0Y_1r_1 \right] \left[ r_1Y_2r_2 \right] \cdots\left[ r_{k-1}Y_k p
      \right] $ for all states $r_1,r_2,\ldots,r_k\in Q$. Let
      $x=\omega_1\omega_2\cdots\omega_k$ where each $\omega_i$ is the input consumed to
      `pop' $Y_i$ off the stack. We have the sequences of less than $n$ moves $\left( r_{i-1},\omega_i,Y_i \right)
      \vdash{*}\left( r_i,\epsilon,\epsilon \right)$ and hence $\left[
      r_{i-1}Y_ir_i \right] \Rightarrow
      ^{*}\omega_i$ by inductive hypothesis.
      \[\left[ qXr_k \right] \Rightarrow a\left[ r_0Y_1r_1 \right] \left[
      r_1Y_1r_2 \right] \cdots\left[ r_{k-1}Y_kp \right] \Rightarrow
    ^{*}aw_1\left[ r_1Y_2r_2 \right] \left[ r_2Y_3r_3 \right] \cdots\left[
  r_{k-1}Y_kp \right] \Rightarrow ^{*}aw_1w_2\cdots w_k.\]
  \end{enumerate}
  We have $\left[ qXp \right]\Rightarrow ^{*}\omega$.\\
  
  ($\impliedby$) Suppose $\left[ qXp \right] \Rightarrow ^{*}\omega$. Inducting
  on the number of steps in the derivation, we have:
  \begin{enumerate}
    \item \textit{Basis Step:} With one step, $\left[ qXp \right] \to \omega$
      is a production $\iff$ there is a transition of $P$ in which $X$ is
      popped and we transit states from $q$ to $p$. $\left( p,\epsilon \right)
      \in \delta\left( q,a,X \right) $ and $a=\omega$, hence $\left( q,\omega,X
      \right) \vdash \left( p,\epsilon,\epsilon \right) $.
    \item \textit{Inductive Step:} Suppose $\left[ qXp \right] \Rightarrow
      ^{*}\omega$ by $n>1$ steps. The derivation must be of the form
      \[\left[ qXp \right] \Rightarrow a\left[ r_0Y_1r_1 \right] \left[
      r_1Y_2r_2 \right] \cdots\left[ r_{k-1}Y_kp \right]\Rightarrow
    ^{*}\omega.\]
    We have $\left( r_0,Y_1Y_2\cdots Y_k \right) \in \delta\left( q,a,X \right) $.
    Let $\omega=a\omega_1\omega_2\cdots\omega_k$ such that $\left[
    r_{i-1}Y_ir_i \right] \Rightarrow ^{*}\omega_i$. By inductive hypothesis,
    $\left( r_{i-1},w_i,Y_i \right) \vdash^{*}\left( r_i,\epsilon,\epsilon
    \right) $ for each $i$. This gives us
    \[\left( q,a\omega_1\omega_2\cdots\omega_k,X \right) \vdash\left(
    r_0,\omega_1\omega_2\cdots\omega_k,Y_1Y_2\cdots Y_k \right)\vdash
  ^{*}\left( r_1,\omega_2\omega_3\cdots\omega_k,Y_2Y_3\cdots Y_k \right)\vdash
^{*}\cdots\vdash ^{*}\left( p,\epsilon,\epsilon \right).\] 
  \end{enumerate}
  Since $S\Rightarrow ^{*}\omega\iff\left[ q_0Z_0p \right]
  \Rightarrow ^{*}\omega$ for some state $p$ (by construction), and each
  $\left[ q_0Z_0p \right]\Rightarrow ^{*}\omega\iff\left( q_0,\omega,Z_0
  \right) \vdash^{*}\left( p,\epsilon,\epsilon \right) $, we have $L(G)=L(P)$
  with acceptance by empty stack.
\end{proof}

\subsection{Context Free Languages}
To show that a given language is non context-free, we may use the pumping
lemma. Intuitively, in any sufficiently long string, it is possible to find two
short nearby substrings that we can ``pump'' in tandem, while keeping the
resulting string in $L$.
\begin{lemma}
  Let $L$ be a CFL\@. Then $\exists n=f(L)$ such that $\forall z\in L$,
  $\left| z \right| \ge n$, $\exists u,v,w,x,y\in \Sigma^{*},z=uvwxy$ such that:
  \begin{enumerate}
    \item $\left| vwx \right| \le n$. That is, the two strings to be ``pumped''
      is not too far from the other.
    \item $vx\neq\epsilon$. At least one of the strings to be ``pumped'' cannot be empty.
    \item $\forall i\ge 0$, $uv^{i}wx^{i}y\in L$
  \end{enumerate}
\end{lemma}
The contrapositive of the lemma is used to prove that a language is
not context-free. That is, $\forall n\in \mathbb{N}$, $\exists z\in L$, such
that $\left| z \right| \ge n$, $\forall u,v,w,x,y\in \Sigma^{*}$, where
$z=uvwxy$, $\left| vwx \right| \le n$, $vx\neq\epsilon$, we have
\[\exists i\in \mathbb{N}\left( uv^{i}wx^{i}y\notin L \right)\]

\begin{proof}
  Suppose $L$ is context-free. Without loss of generality, assume
  $L\neq\emptyset $, $L\neq\left\{ \epsilon \right\} $ and a CNF grammar
  $G=\left( V,T,P,S \right) $ for $L-\left\{ \epsilon \right\} $. Let $m=\left|
  V \right| $ and $n=2^{m}$. Consider a string $z\in L$ such that $\left| z
  \right| \ge n$.  Its parse tree must have a path from the
  root to a leaf of length at least $m+1$. By Pigeonhole Principle, there must
  exist a repeated nonterminal $A$.\\

  We have the derivation of the form, where $z=uvwxy$ for some $u,v,w,x,y\in \Sigma^{*}$,
  \[S\Rightarrow ^{*}uAy\Rightarrow ^{*}uvAxy\Rightarrow ^{*}uvwxy.\]
  \begin{figure}[h!]
    \centering
    \includegraphics[width=0.5\textwidth]{cfl_pumping_lemma}
  \end{figure}

  Additionally, we have $A\Rightarrow ^{*}vAx$ and $A\Rightarrow ^{*}w$. For
  any $i\in \mathbb{N}$, we thus may have the derivation
  \[S\Rightarrow ^{*}uAy\Rightarrow ^{*}uv^{i}Ax^{i}y\Rightarrow ^{*}uv^{i}wx^{i}y.\]
  With the assumption that the grammar is of CNF, since the repeated terminals
  are maximally $m$ production steps apart, we have the length $\left|  vwx
  \right| \le n$, and since the derivation $A\Rightarrow ^{*}vAx$ uses at least
  one production step, $vx\neq \epsilon$.\\

  For the excluded cases when constructing the CNF grammar, the statement of
  the pumping lemma is vacuously true for $L=\emptyset $ and we choose $n$ such
  that $z\neq\epsilon$ for $L=\left\{ \epsilon \right\} $.
  \qed%
\end{proof}

By using the pumping lemma, we can construct a proof of the following structure to show a specified language is non context-free:

\begin{prop}
  $L=\left\{ a^{m}b^{m}c^{m} | m\ge 1\right\} $ is not context-free.
\end{prop}
\begin{proof}
  Suppose $L$ is regular. Let $n$ be as in the Pumping Lemma. Let
  $z=a^{n}b^{n}c^{n}$. Let $z=uvwxy$ as in the pumping lemma. Since $\left| vwx
  \right| \le n$, $vwx$ cannot contain both $a$ and $c$. \\

  Suppose $vwx$ does not contain $a$'s. By the pumping lemma, $uv^{2}wx^{2}y\in
  L$. But $uv^{2}wx^{2}y$ contains $n$ $a$'s, and $\left| uv^{2}wx^{2}y \right|
  >3n$, hence $uv^{2}wx^{2}y\not\in L$.\\

  Suppose $vwx$ does not contain $c$'s. Similarly, By the pumping lemma,
  $uv^{2}wx^{2}y\in L$. But $uv^{2}wx^{2}y$ contains $n$ $c$'s, and $\left|
  uv^{2}wx^{2}y \right| >3n$, hence $uv^{2}wx^{2}y\not\in L$.\\

  Hence $L$ is non context-free.\qed%
\end{proof}

\begin{prop}
  $L=\left\{ \alpha\alpha|\alpha\in \left\{ a,b \right\} ^{*} \right\} $ is not
  context-free.
\end{prop}

\begin{lemma}[Ogden's Lemma]
  Let $L$ be a CFL. Then $\exists n$ such that $\forall z\in L$, $\left| z
  \right| \ge n$, $\exists u,v,w,x,y\in \Sigma^{*}$, $z=uvwxy$, if $\ge n$ positions
  in $z$ are marked as distinguished,
  \begin{itemize}
    \item $vwx$ has at most $n$ distinguished positions.
    \item $vx$ has at least $1$ distinguished position.
    \item $\forall i\ge 0$, $uv^{i}wx^{i}y\in L$.
  \end{itemize}
\end{lemma}
Ogden's Lemma provides an stronger form of the pumping lemma that can be used
to prove non-context free-ness.

\subsubsection{Closure of Context-Free Languages}
We have, similarly to regular languages, certain operations that, when applied to CFLs,
the resulting language is also a CFL\@.

\begin{definition}[Substitution]
  Suppose terminals $i$ and CFLs $L_i$. Consider a mapping
  $s:\Sigma\to\mathcal{P}\left( \Gamma^{*} \right) $ such that $s(i)=L_i$, that
  is further defined for strings by principle of induction:
  \begin{enumerate}
    \item \textit{Base Case:} $s(\epsilon)=\left\{ \epsilon \right\}$
    \item \textit{Inductive Case:} $s(\omega a)=s(\omega)\cdot s(a), \forall a\in \Sigma,\omega\in \Sigma^{*}$
  \end{enumerate}
  $s$ is called a \textit{substitution}. Intuitively, we replace each symbol in the
  strings of one language by an entire language.
\end{definition}
\begin{theorem}
  Given a CFL $L$ over $\Sigma$ and $s$ a substitution on $\Sigma$ such that
  $s\left( a \right) $ is a CFL for each $a\in \Sigma$, $s\left( L \right) $ is
  a CFL.
\end{theorem}
The principle of the proof is to consider replacing each terminal in the
original grammar with a nonterminal that acts as the start symbol of the
substituted CFL. We then have a resultant CFG that generates $s\left( L \right)
$.\\

Given CFLs $L,L_1,L_2$ and regular language $R$:
 \begin{itemize}
  \item $L_1\cup L_2$ is context-free.
  \item $L_1\cdot L_2$ is context-free.
  \item $L^{*}$ is context-free.
  \item $L^{+}$ is context-free.
  \item $h\left( L \right) $ is context-free, for some homomorphism $h$.
  \item $L^{R}$ is context-free. Consider the reverse of all productions.
  \item $L\cap R$ is context-free. Consider the PDA for $L$, and DFA for
    $R$. We can simulate the DFA with the analogue for parallelisation of
    DFAs.
  \item $h^{-1}\left( L \right) $ is context-free, for some homomorphism $h$.
\end{itemize}

\begin{remark}
  $\bar{L}$ and $L_1-L_2$ are not necessarily context-free.
\end{remark}

\begin{remark}
  CFLs are not closed under intersection. Consider the CFLs $L_1=\left\{
  0^{n}1^{n}2^{i}|n\ge 1,i\ge 1 \right\} $ and $L_2=\left\{
0^{i}1^{n}2^{n}|n\ge 1,i\ge 1 \right\} $. We have the intersection $L=\left\{
0^{n}1^{n}2^{n}|n\ge 1 \right\} $ which is not context-free.
\end{remark}

\subsubsection{Decision Problems of CFLs}
We may examine some decision problems of CFLs. However, in CFLs, we have to
take into account the complexities of converting between the various forms of
representation, that is required for their decision problems.\\

We have the linear time conversions from a CFG to a PDA, from a PDA accepting
by final state to a PDA accepting by empty stack and vice-versa. On the other hand, the
conversion from a PDA to a CFG is requires an upper bound of $O(n^{3})$ on the
constructed grammar's length and hence running time. Additionally, to convert a
general CFG to its CNF requires $O(n^{2})$ for the length of the resultant
grammar, and running time to construct the unit pairs and eliminate unit
productions.\\

Given these conversions, we are still highly limited in what we may decide
about a CFL, many of which is undecideable as shown in a later chapter.\\

We may check the \textbf{emptiness of CFLs} by running the algorithm to find
all generating symbols in its CFG and checking whether the start symbol $S$ is
generating. That is,
\begin{gather*}
  L=\emptyset \iff  S\text{ is not generating}
\end{gather*}

We may decide the \textbf{membership of a string} $\omega=a_1a_2\cdots a_{n}$
in a CFL $L$. This can be done with dynamic programming, using the CYK
algorithm using a CNF grammar for the language. We construct a triangular table
where each entry $X_{ij}$ is the set of variables $A$ such that $A\Rightarrow
^{*}a_ia_{i+1}\cdots a_j$.
\begin{figure}[h!]
  \centering
  \includegraphics[width=0.3\textwidth]{cyk_algo}
  \captionsetup{labelformat=empty}
  \caption{Table constructed by CYK Algorithm for $n=5$}
\end{figure}\\
We compute the rows bottom-up, with the following inductive algorithm:
\begin{enumerate}
  \item \textit{Base Case:} Each $X_{ii}$ is the set of nonterminals which generate $a_i$
  \item  \textit{Inductive Step:} $X_{i,j}$ contains all $A$ such that $A\to BC
    $, where $B\in X_{i,k}$ and $C\in X_{k+1,j}$ for some $i\le k<j$.
\end{enumerate}
We then have the membership of the string if the start symbol generates the
full length of the string. That is,
\[\omega\in L\iff S\in X_{1,n}\]
\begin{remark}
  The running time of the CYK algorithm is $O(n^{3})$, in comparison to naive
  recursive searches for membership of strings which require exponential time.
\end{remark}

\begin{remark}
  The following CFL decision problems are undecideable:
  \begin{itemize}
    \item Is a given CFG $G$ ambiguous?
    \item Is a given CFL inherently ambiguous?
    \item Is the intersection of two CFLs empty?
    \item Are two CFLs the same?
    \item Is a given CFL equal to $\Sigma^{*}$?
    \item Is a given CFL a regular language?
  \end{itemize}
\end{remark}

\subsection{Deterministic PDA}
PDAs are nondeterministic by definition. However, parsers generally behave
deterministically, and not unlike a deterministic PDA\@. We investigate, in this
section, the behaviour and types of languages accepted by deterministic PDAs.

\begin{definition}[DPDA]
  An PDA $P=\left( Q,\Sigma,\Gamma,\delta,q_0,Z_0,F \right) $ is a \textit{DPDA} if:
  \begin{enumerate}
    \item $\forall q\in Q$, $a\in \left\{ \epsilon \right\} \cup \Sigma$, and
      $X\in \Gamma$,$\qquad\left| \delta\left( q,a,X \right) \right| \le 1$.
      There is only one possible move, if any, for any input and ID.
    \item If $\delta\left( q,\epsilon,X \right)\neq \emptyset\implies\forall
      a\in \Sigma, \delta\left( q,a,X \right)=\emptyset $. That is, we do not
      have a choice between taking a input symbol or making a move on
      $\epsilon$.
  \end{enumerate}
\end{definition}

\subsubsection{Languages Accepted by DPDAs}
We define the languages accepted by DPDAs similarly to that of PDAs.

\begin{theorem}[Languages Accepted by Final State]
  The languages accepted by DPDAs by final state properly include the regular
  languages, but are properly included in the CFLs.
\end{theorem}
On the other hand, DPDAs that accept by empty stack are highly limited, unable
to accept even certain finite languages such as $\left\{ a,aa \right\}$.
\begin{definition}[Prefix Property]
  A language $L$ has the \textit{prefix property} if $\forall \left( x,y
  \right) \in L^{2}$, $x\text{ is not a prefix of }y$.
\end{definition}
\begin{theorem}[Languages Accepted by Empty Stack]
  A language $L$ is $L(P)$ by empty stack for some DPDA $P\iff$ $L$ has the
  prefix property $L$ is $L(P')$ by final state for some DPDA $P$.
\end{theorem}

\subsubsection{Unambiguous Grammars}
\begin{theorem}
  If $L=L(P)$ by acceptance through empty stack for some DPDA $P$, $L$ has an
  unambiguous context-free grammar.
\end{theorem}
The typical conversion of CFG from PDA (with nonterminals encoding state
transition and symbol popped) yields an unambiguous context-free grammar.
Intuitively, this is because there is a unique sequence of moves in the DPDA,
and hence there is only one choice of production every step of the leftmost
derivation of each string. There is also only one final state that can be
reached due to the determinism, and hence only one of the initial branches of
$S$ will be consistent with the PDA\@.

\begin{theorem}
  If $L=L(P)$ by acceptance through final state for some DPDA  $P$, $L$ has an
  umabiguous context-free grammar.
\end{theorem}
This can be shown by adding an end marker $\$\not\in\Sigma$ to every string in
$L$, then creating the context free grammar as though for acceptance through
empty stack, and then introducing a production $\$\to\epsilon$.
